\documentclass[numbers=noenddot,bibliography=totoc,12pt,sfdefaults=false]{scrartcl}
\usepackage{geometry}
\geometry{a4paper}
\geometry{twoside=false, vmargin=1cm, hmargin=1in, centering, columnsep=2pc, heightrounded, includeheadfoot}
%%%%%%%%%%%%%%%%%%%%%%%%%%%%%%%%%%%%%%%%%%%%%%%%%%%%%%%%%%%%%%%%%%%%%%%%%%%%

%\newcommand{\lang}{ngerman} % english or ngerman
\newcommand{\lang}{ngerman} % english or ngerman
\newcommand{\name}{Marco Garcia Jung}
\newcommand{\street}{Klenzestraße 25}
\newcommand{\plzCity}{TODO PLZ und Ort}
\newcommand{\matnr}{2437251}

\newcommand{\submittedDate}{16. Februar 2027}

\newcommand{\examiner}{Prof. Dr. Stefan Schönig}
\newcommand{\supervisor}{Maximilian Metzner} % TODO prüfen

\newcommand{\thesistypeDe}{Bachelorarbeit}
\newcommand{\thesistypeEn}{Bachelor's Thesis}

\newcommand{\thesisTitle}{Tabular Foundation Models bei fehlenden Merkmalswerten: Ein Vergleich mit Gradient Boosting bei Klassifikationsaufgaben}
%%%%%%%%%%%%%%%%%%%%%%%%%%%%%%%%%%%%%%%%%%%%%%%%%%%%%%%%%%%%%%%%%%%%%%%%%%%%

\newcommand{\thesistype}{not set}
\newcommand{\abbreviationname}{not set}
\usepackage[\lang]{babel}

\addto\captionsenglish{
    \renewcommand{\listfigurename}{List of Figures}
    \renewcommand{\listtablename}{List of Tables}
    \renewcommand{\listoflistingscaption}{List of Code Examples}
    \renewcommand{\abbreviationname}{Abbreviations and Acronyms}
    \crefname{listing}{code example}{code examples}
    \Crefname{listing}{Code Example}{Code Examples}
    \captionsetup[listing]{name=Code}
    \sisetup{locale=US}
    \renewcommand{\thesistype}{\thesistypeEn}
}

\addto\captionsngerman{
    \renewcommand{\listfigurename}{Abbildungsverzeichnis}
    \renewcommand{\listtablename}{Tabellenverzeichnis}
    \renewcommand{\listoflistingscaption}{Codebeispiele}
    \renewcommand{\abbreviationname}{Abkürzungen}
    \crefname{listing}{Codebeispiel}{Codebeispiele}
    \Crefname{listing}{Codebeispiel}{Codebeispiele}
    \captionsetup[listing]{name=Codebeispiel}
    \sisetup{locale=DE}
    \renewcommand{\thesistype}{\thesistypeDe}
}

\sloppy

\usepackage{minted}
\setminted{frame=lines, framesep=2mm, linenos, breaklines=true, fontsize=\footnotesize,style=manni,autogobble}

\usepackage{amsmath}

\usepackage[T1]{fontenc}
\usepackage[p,osf,cochineal,swashQ,vvarbb]{newtx}
\usepackage[book,scale=0.82]{FiraSans}
\usepackage[cal=boondoxo]{mathalfa}
\usepackage[scale=0.82]{FiraMono}

\usepackage{setspace}
\setstretch{1}

\usepackage[hang,flushmargin]{footmisc}
\usepackage[svgnames,dvipsnames]{xcolor}
\usepackage{iftex}
\usepackage{microtype}
\usepackage{hyperref}
\hypersetup{%
  colorlinks,
  bookmarksnumbered,
  pdfencoding=unicode,
  allcolors=DarkSlateGrey
}

\setkomafont{section}{\rmfamily\fontsize{13pt}{15pt}\bfseries}
\setkomafont{subsection}{\rmfamily\fontsize{11pt}{13pt}\bfseries\selectfont}
\setkomafont{subsubsection}{\rmfamily\fontsize{10.5pt}{12pt}\fontseries{b}\fontshape{it}\selectfont}
\setkomafont{paragraph}{\rmfamily\fontsize{11pt}{10pt}\itshape\selectfont}
\setkomafont{subparagraph}{\rmfamily\fontsize{11pt}{10pt}\selectfont}

\colorlet{scolor}{black}
\colorlet{hscolor}{DarkSlateGrey}
\usepackage{booktabs}
\usepackage[backend=biber, style=lncs, giveninits,sorting=nyt]{biblatex}
\addbibresource{../references.bib}

\usepackage{siunitx}
\sisetup{per-mode = fraction, prefixes-as-symbols = false, list-units = brackets, range-units = brackets, multi-part-units = brackets }

\usepackage{graphicx}
\graphicspath{{figures/}}

\usepackage{enumitem}
\usepackage{cleveref}
\usepackage{tabularx}
\usepackage[printonlyused]{acronym}

\usepackage{tabularray}
\UseTblrLibrary{booktabs}
\DefTblrTemplate{firsthead, middlehead,lasthead}{default}{}

\usepackage{csquotes}
\MakeOuterQuote{"}

\usepackage{tikz}
\usetikzlibrary{3d,mindmap}

\usepackage{caption}
\captionsetup{font={small}, labelfont={bf,small}, singlelinecheck=false, labelsep=newline, format=plain}

\usepackage[colorinlistoftodos,prependcaption,textsize=tiny]{todonotes}

\usepackage[defaultlines=3,all]{nowidow}

% Misc
\usepackage{ifthen}
\usepackage{iflang}

\usepackage{scrlayer-scrpage}
\automark*[section]{}
\renewcommand*{\partmark}[1]{}
\setkomafont{pageheadfoot}{\normalcolor\footnotesize\scshape}
\setkomafont{pagenumber}{\normalfont\normalsize}
\clearpairofpagestyles
%\ihead{\headmark}
\ohead*{\normalcolor\footnotesize\pagemark}
\ifoot*{\normalfont \thesistype}
\ofoot*{\normalfont \name, \the\year{}}
\KOMAoptions{headsepline=.25pt}
\pagestyle{headings}
\usepackage{scrhack}

%%
%% end of the preamble, start of the body of the document source.
\begin{document}

    \thispagestyle{empty}
\begin{titlepage}

\renewcommand{\thepage}{}

\begin{center}

{Universität Regensburg\\}
{Fakultät für Wirtschaftswissenschaften\\}
{Lehrstuhl für Prozessbasierte Informationssysteme\\}


\vspace*{10mm}

\Large{\textbf{\thesisTitle}}

\vspace*{10mm}
\includegraphics[width=0.55\textwidth]{ur-logo-bildmarke-grau}
\vspace*{4mm}

\Large{\thesistypeDe}

\vspace*{5mm}


\end{center}
\begin{center}
%\ifthenelse{\equal{\documenttype}{Bachelorarbeit} \or \equal{\documenttype}{Bachelorthesis}}{Zur Erlangung des akademischen Grades "Bachelor of Science (M.Sc.)" im Studiengang Wirtschaftsinformatik an der Fakultät für Informatik und Data Science der\\Universität Regensburg gemäß §45 Masterprüfungsordnung vom 28.09.2015}{}

%\ifthenelse{\equal{\documenttype}{Masterarbeit} \or \equal{\documenttype}{Masterthesis}}{Zur Erlangung des akademischen Grades "Master of Science (M.Sc.)" im Studiengang Wirtschaftsinformatik an der Fakultät für Informatik und Data Science der\\Universität Regensburg gemäß §45 Masterprüfungsordnung vom 28.09.2015}{}

%Zur Erlangung des akademischen Grades "Master of Science (M.Sc.)" im Studiengang Wirtschaftsinformatik an der Fakultät für Informatik und Data Science der\\ Universität Regensburg gemäß §45 Masterprüfungsordnung vom 28.09.2015

\vspace*{10mm}

\Large{Eingereicht bei: \examiner\\}
\Large{Betreuung: \supervisor\\}


\end{center}

\vfill

\begin{center}
\large{Eingereicht am \submittedDate\\}
\end{center}
\vspace*{0.6cm}
\begin{flushleft}
\textit{Eingereicht von}\\
\vspace*{5pt}
\name\\
\street\\
\plzCity\\
\vspace*{5pt}
Matrikelnummer: \matnr \\

\end{flushleft}
\end{titlepage}

\newpage


% Abstract wird zuletzt geschrieben
    {\thispagestyle{empty}
    \subsection*{\abstractname}
    }

    \newpage

    \pagenumbering{roman}
    \setcounter{page}{1}
    \tableofcontents
    \newpage

    \addcontentsline{toc}{section}{\listfigurename}
    \listoffigures
    \newpage

    \addcontentsline{toc}{section}{\listtablename}
    \listoftables
    \newpage

% Nur einkommentieren, falls Codebeispiele in der Arbeit stehen
%    \addcontentsline{toc}{section}{\listoflistingscaption}
%    \listoflistings
%    \newpage

% Es erscheinen nur Abkürzungen, die im Text mit \ac{...} verwendet werden
    \section*{\abbreviationname}
    \addcontentsline{toc}{section}{\abbreviationname}
    \begin{acronym}[MMMMM]
        \acro{BA}{Balanced Accuracy}
        \acro{ICL}{In-Context Learning}
        \acro{LLM}{Large Language Model}
        \acro{MAR}{Missing At Random}
        \acro{MCAR}{Missing Completely At Random}
        \acro{MNAR}{Missing Not At Random}
        \acro{PFN}{Prior-Data Fitted Network}
        \acro{TFM}{Tabular Foundation Model}
    \end{acronym}
    \newpage
    \pagenumbering{arabic}
    \setcounter{page}{1}

%%%%%%%%%%%%%%%%%%%%%%%%%%%%%%%%%%%%%%%%%%%%%%%%%%%%%%%%%%%%%%%%%%%%

    % Seitenbudget: ca. 4 Seiten. Wird zuletzt geschrieben.
\section{Einleitung}\label{sec:einleitung}

\subsection{KI, LLMs und der Übergang zu Tabular Foundation Models}\label{sec:einleitung-kontext}

\subsection{Problemstellung: bestehende Modellvergleiche und der Fokus auf fehlende Merkmalswerte}\label{sec:einleitung-problem}

\subsection{Zielsetzung, Beitrag und Forschungsfragen}\label{sec:einleitung-ziel}

\subsection{Abgrenzung und Aufbau}\label{sec:einleitung-aufbau}

    % Seitenbudget: ca. 7 Seiten. Nur, was FF1 und FF2 brauchen.
\section{Grundlagen}\label{sec:grundlagen}

\subsection{Klassifikation auf tabellarischen Daten und Datenbesonderheiten}\label{sec:grundlagen-tabellen}

\subsection{Transformer, Foundation Models und In-Context Learning}\label{sec:grundlagen-icl}

\subsection{Bayesianische Vorhersage und Prior-Data Fitted Networks}\label{sec:grundlagen-pfn}

\subsection{Tabular Foundation Models am Beispiel TabPFN}\label{sec:grundlagen-tabpfn}

\subsection{Gradient Boosting am Beispiel LightGBM}\label{sec:grundlagen-lightgbm}

\subsection{Fehlende Werte: MCAR, MAR, MNAR, Imputation und Missing-Indikatoren}\label{sec:grundlagen-missing}

\subsection{Evaluationsmetriken}\label{sec:grundlagen-metriken}

    % Seitenbudget: ca. 6 Seiten. Beantwortet FF1.
% Quellen: literature/notes/Suchprotokoll.md und literature/notes/Konzeptmatrix.md
\section{Literaturanalyse zu fehlenden Merkmalswerten}\label{sec:literatur}

\subsection{Recherche, Auswahl und Konzeptmatrix}\label{sec:literatur-vorgehen}

\subsection{Bestehende Leistungsvergleiche zwischen Tabular Foundation Models und Gradient Boosting}\label{sec:literatur-vergleiche}

\subsection{Verarbeitung fehlender Werte in beiden Modellansätzen}\label{sec:literatur-verarbeitung}

\subsection{Empirische Befunde zu Fehlmechanismen, Fehlraten und Behandlungen}\label{sec:literatur-befunde}

\subsection{Bewertung der Evaluationsdesigns und ihrer Aussagekraft}\label{sec:literatur-bewertung}

\subsection{Abgeleiteter Untersuchungsbedarf und Anforderungen an das Experiment}\label{sec:literatur-bedarf}

    % Seitenbudget: ca. 6 Seiten. Design für FF2.
\section{Methodik der empirischen Untersuchung}\label{sec:methodik}

\subsection{Forschungsdesign und Ableitung aus der Literaturanalyse}\label{sec:methodik-design}

\subsection{Verarbeitungsprinzipien und Auswahl der Vertreter: TabPFN-3.5 und LightGBM}\label{sec:methodik-vertreter}

\subsection{Datensatzauswahl, gruppengetrennte Splits identischer Eingaben und Vorverarbeitung}\label{sec:methodik-daten}

\subsection{Erzeugung fehlender Werte: Fehlmechanismen und Fehlraten}\label{sec:methodik-masken}

\subsection{Interne Verarbeitung, Median-Imputation und Indikator-Zusatzanalyse}\label{sec:methodik-behandlungen}

\subsection{Konfigurationen und Tuning-Sensitivität}\label{sec:methodik-tuning}

\subsection{Wiederholungen, Metriken und gepaarte Auswertung}\label{sec:methodik-auswertung}

\subsection{Pilotlauf, fixierte Software- und Checkpointversionen und Reproduzierbarkeit}\label{sec:methodik-repro}

    % Seitenbudget: ca. 9 Seiten. Kern der Arbeit, beantwortet FF2.
\section{Ergebnisse}\label{sec:ergebnisse}

\subsection{Datensatzübersicht und Ausgangsleistung}\label{sec:ergebnisse-referenz}

\subsection{Leistungskurven nach Fehlrate, Behandlung und Fehlmechanismus}\label{sec:ergebnisse-kurven}

\subsection{Veränderung des gepaarten Leistungsunterschieds und Indikator-Pipeline}\label{sec:ergebnisse-abstand}

\subsection{Sensitivität gegenüber erneuter Abstimmung des Gradient Boosting}\label{sec:ergebnisse-tuning}

\subsection{Streuung und Konsistenz über Datensätze}\label{sec:ergebnisse-streuung}

    % Seitenbudget: ca. 6 Seiten.
\section{Diskussion}\label{sec:diskussion}

\subsection{Beantwortung von FF1}\label{sec:diskussion-ff1}

\subsection{Beantwortung von FF2}\label{sec:diskussion-ff2}

\subsection{Synthese und Beantwortung der übergeordneten Forschungsfrage}\label{sec:diskussion-synthese}

\subsection{Praktische Bedeutung, Limitationen und Übertragbarkeit}\label{sec:diskussion-limitationen}

    % Seitenbudget: ca. 2 Seiten.
\section{Fazit und Ausblick}\label{sec:fazit}


%%%%%%%%%%%%%%%%%%%%%%%%%%%%%%%%%%%%%%%%%%%%%%%%%%%%%%%%%%%%%%%%%%%%
    \clearpage

    \printbibliography

    \clearpage
    \appendix

    % Der Anhang zählt nicht zum Seitenbudget.
\section{Dokumentation der Literaturrecherche}\label{sec:anhang-recherche}

\section{Ergänzende Ergebnistabellen}\label{sec:anhang-ergebnisse}

\section{Software, Checkpoint und Seeds}\label{sec:anhang-repro}


%%%%%%%%%%%%%%%%%%%%%%%%%%%%%%%%%%%%%%%%%%%%%%%%%%%%%%%%%%%%%%%%%%%%

\newpage

\subsection*{Angaben zur Verwendung KI-basierter Hilfsmittel}

Ich versichere, dass ich keine KI-basierten Tools verwendet habe, deren Nutzung der Prüfer / die Prüferin explizit schriftlich ausgeschlossen hat. Sämtliche KI-basierte Hilfsmittel sind mit Produktnamen und Erläuterung der Verwendung in folgender Tabelle angegeben. Ich bin mir bewusst, dass die Verwendung von Texten oder anderen Inhalten und Produkten, die durch KI-basierte Tools generiert wurden, keine Garantie für deren Qualität darstellt. Ich verantworte die Übernahme jeglicher von mir verwendeter maschinell generierter Passagen vollumfänglich selbst und trage die Verantwortung für eventuell durch die KI generierte fehlerhafte oder verzerrte Inhalte, fehlerhafte Referenzen, Verstöße gegen das Datenschutz- und Urheberrecht oder Plagiate. Ich versichere zudem, dass in der vorliegenden Arbeit mein gestalterischer Einfluss überwiegt.

\vspace{2cm}
\begin{flushright}
    Regensburg, den \submittedDate
\end{flushright}

% Laufend pflegen: jedes Tool mit Einsatzform, Grund und betroffenen Teilen
{
\scriptsize
\begin{longtblr}[entry=none]{
  width = \linewidth,
  colspec = {Q[l,-1]X[l]X[l]Q[l]X[l]},
  rowhead =
}
\toprule
\SetCell{font=\bfseries} KI-Tool & \SetCell{font=\bfseries} Einsatzform & \SetCell{font=\bfseries} Grund & \SetCell{font=\bfseries} Betroffene Teile & \SetCell{font=\bfseries} Anmerkungen \\
\midrule

GitHub Copilot & Einrichtung von Projektstruktur und LaTeX-Gerüst, Arbeitsplanung & Organisation des Arbeitsprozesses & Gesamte Arbeit (Struktur) & Kein Fließtext übernommen \\

\bottomrule
\end{longtblr}
}


\newpage
\pagestyle{empty}

\subsection*{Erklärung an Eides statt}
\todo[inline]{Geltende Prüfungsordnung prüfen: unten steht die Fassung BPO 2021; für PO 2015 lautet der Verweis §13 Abs. 3 PO 2015.}
Ich habe die vorliegende Arbeit selbstständig verfasst und keine anderen als die angegebenen Quellen und Hilfsmittel benutzt. Die Arbeit wurde bisher an keiner anderen Hochschule zur Erlangung eines akademischen Grades eingereicht. Die vorgelegten Druckexemplare und die dem Prüfer / der Prüferin zur Verfügung gestellte elektronische Version (PDF-Datei) der Arbeit sind identisch. Von den in §27 Abs. 6 BPO 2021 vorgesehenen Rechtsfolgen habe ich Kenntnis.

\vspace{4cm}
\begin{flushright}
    Regensburg, den \submittedDate
\end{flushright}

\end{document}
